% ============================================================
%  Søren Kierkegaard, "Ved en Grav"
%  (Tre Taler ved tænkte Leiligheder, Kjøbenhavn 1845 --- tredje tale)
%
%  TRANSCRIPTION (Danish)
%  Source: Søren Kierkegaards Skrifter (SKS), reading text ("Hovedtekst"),
%  Det Kgl. Biblioteks tekstportal,
%  https://tekster.kb.dk/text/sks-ttl-txt-root  (SKS bind 5, s. 442--469).
%  Original 1845 orthography retained. Editorial page numbers and
%  apparatus are NOT reproduced; only Kierkegaard's own text is given.
% ============================================================
\documentclass[12pt,a4paper]{article}
\usepackage[utf8]{inputenc}
\usepackage[T1]{fontenc}
\usepackage{libertinus}
\usepackage{libertinust1math}
\usepackage[danish]{babel}
\usepackage[protrusion=true,expansion=true]{microtype}
\usepackage{setspace}
\usepackage[margin=1.2in]{geometry}
\usepackage{fancyhdr}
\usepackage{hyperref}

\hypersetup{
  pdftitle={Ved en Grav},
  pdfauthor={Søren Kierkegaard},
  hidelinks
}

\pagestyle{fancy}
\fancyhf{}
\rhead{\emph{Tre Taler ved tænkte Leiligheder} (1845)}
\cfoot{\thepage}

\setstretch{1.3}

\title{\textbf{Ved en Grav}}
\author{S.\ Kierkegaard}
\date{\emph{Tre Taler ved tænkte Leiligheder}, Kjøbenhavn 1845 --- tredje tale\\[2pt]
  {\small SKS 5, s.\ 442--469}}

\begin{document}
\maketitle
\thispagestyle{fancy}
\bigskip

\noindent Saa er det da forbi! – Og naar nu Den, som her først traadte hen til Graven, fordi han er den Nærmeste, naar han efter Talens korte Øieblik er bleven den Sidste ved Graven, ak, fordi han er den Nærmeste: saa er det forbi. Om han vilde blive derude, han erfarer dog ikke hvad den Døde foretager sig, thi den Døde er en stille Mand; om han i sin Bekymring vilde kalde paa hans Navn, om han i sin Sorg vilde sidde lyttende, han erfarer dog Intet, thi i Graven er der Stilhed, og den Døde er en taus Mand; og om han erindrende gik hver Dag til hans Grav, den Døde erindrer ikke ham –.

Thi i Graven er der ingen Erindring, end ikke af Gud. See, dette vidste den Mand, om hvem der nu maa siges, at han ikke mere erindrer Noget, til hvem det nu var for silde at sige det. Men som han vidste det, saa gjorde han derefter, og derfor erindrede han Gud medens han levede. Hans Liv gik hen i hæderlig Ubemærkethed, ikke Mange vare vidende om hans Tilværelse, kun Enkelte blandt de Faa kjendte ham. Han var Borger her af Staden, stræbsom i sin beskedne Gjerning forstyrrede han Ingen ved at svigte Samfundets Forpligtelser, forstyrrede Ingen ved utidig Bekymring om det Hele. Saaledes gik Aar efter Aar, eensformigt, dog ikke tomt; han blev Mand, han blev gammel, han blev bedaget: Gjerningen var og blev den samme, een og samme Syssel i de forskjellige Levealdre. Han efterlader sig en Hustru, glad fordum ved at forenes med ham, nu en gammel Kone, der begræder den Tabte, en ret Enke, der, forladt, har sit Haab til Gud. Han efterlader sig en Søn, der lærte at elske ham og at finde Tilfredshed i Vilkaaret og Faderens Gjerning; engang glad som Barn i Faderens Huus, fandt han det som Yngling aldrig for trangt, nu er det ham et Sørgehuus. – En saadan ubemærket Mands Død spørges ikke vidt, og naar Nogen om føie Tid gaaer forbi det Huus, hvor han boede i Ringhed, og læser hans Navn over Døren, fordi den borgerlige Bedrift fortsættes under hans Navn, da er det jo, som var han ikke død. Som han sov mildt og fredeligt hen, saaledes er i den omgivende Verden hans Død en Bortgang i Stilhed. Brav som Borger, redelig i sin Næringsvei, sparsom i sit Huus, godgjørende efter Evne, deeltagende i Oprigtighed, sin Hustru tro, sin Søn en Fader: see, alt Dette, og al den Sandhed hvormed det her kan siges, spænder ikke Forventningen paa en betydningsfuld Udgang, her er det et Livs Foretagende, paa hvilket en rolig Død blev den skjønne Udgang. – Dog havde han een Gjerning endnu, den blev udført med samme Troskab i Hjertets Eenfoldighed: han erindrede Gud. Han var Mand, gammel, han blev bedaget, saa døde han, men Erindringen om Gud blev den samme, en Veiledning i alt hans Foretagende, en stille Glæde i den fromme Betragtning. Ja, var der slet Ingen, som savnede ham i Døden, ja, var han nu ikke hos Gud, da vilde Gud savne ham i Livet, og vide hans Bopæl og søge ham der, thi den Afdøde vandrede for hans Aasyn og var bedre kjendt af ham end af nogen Anden. Han erindrede Gud, og han blev duelig til sit Arbeide, han erindrede Gud og han blev glad ved sit Arbeide og glad ved Livet, han erindrede Gud, og han blev lykkelig i sit beskedne Hjem med sine Kjære, han forstyrrede Ingen ved Ligegyldighed mod en offentlig Gudsdyrkelse, han forstyrrede Ingen ved ubetimelig Iver, men Guds Huus var ham et andet Hjem – nu er han vandret hjem.

Men i Graven er der ingen Erindring – derfor bliver den tilbage, den bliver hos de Tvende, der vare ham kjære i Livet: de ville erindre ham. Og naar nu Den, som traadte først hen til Graven, fordi han er den Nærmeste, efter Talens korte Øieblik er bleven den Sidste ved Graven, fordi han jo er den Nærmeste, naar han gaaer erindrende herfra, da gaaer han hjem til den sørgende Enke: og Navnet over Døren bliver en Erindring. Saa vil der i nogen Tid nu og da komme en Kjøber, der tilfældigt eller mere deeltagende spørger efter Manden; og naar han hører om hans Død vil Kjøberen sige: naa, er han død. Og naar saa alle de gamle Kjøbere har gjort det eengang, saa har Omgivelsens Liv intet Middel mere for at bevare Erindringen om ham. Men den gamle Enke vil ingen Paamindelse behøve for at erindre, og den stræbsomme Søn vil ikke finde det Sinkende at erindre. Og naar saa Ingen spørger mere om ham, da vil dog Navnet over Døren, naar Huset ikke synligt er et Sørgehuus, naar ogsaa i Huset Sorgen er formildet og det daglige Savn med Trøsten har indøvet Erindringen: da vil Navnet over Døren for de Tvende betyde, at ogsaa de have een Gjerning mere: at erindre den Afdøde.

Nu er Talen forbi. Kun en Handling staaer tilbage: med de tre Spader Jord at indvie den Afdøde, som Alt hvad der er kommen af Jord, til Jord igjen – og saa er det forbi.

Den umyndige Tale kan ikke saaledes gjøre Alvor deraf, ingen Død venter paa den, for at Alt kan være forbi. Men derfor kan Du jo gjerne agte paa Talen, min Tilhører. Thi Døden selv har jo sin Alvor; det Alvorlige ligger ikke i Begivenheden, ikke i det Udvortes: at der nu igjen er en Mand død, saa lidet som Alvorens Forskjel ligger i, at der var mange Karreter; ja saa lidet som hiin formildede Stemning, der kun vil tale godt om de Døde, er Alvor eller i fjerneste Maade kunde tilfredsstille Den, der alvorligen betænkte sin egen Død. Døden kan netop lære, at Alvor ligger i det Indvortes, i Tanken, lære, at det kun er et Sandsebedrag, naar der letsindigt eller tungsindigt sees paa det Udvortes, eller naar Betragteren dybsindigt over Dødens Tanke glemmer at tænke og betænke sin egen Død. Vil man nævne ret en Gjenstand for Alvor, da nævner man Døden, og »Dødens alvorlige Tanke«; og dog er det jo som laae der en Spøg til Grund for Døden, og denne Spøg, forskjelliggjort i Stemnings og Udtryks Forskjellighed, er det Væsentlige i enhver Betragtning af Døden, hvor Betragteren ikke selv bliver under fire Øine med Døden og tænker sig selv sammen med Døden. En Hedning har allerede sagt, at Døden skal man ikke frygte, »thi naar den er, er jeg ikke, og naar jeg er, er den ikke.« Dette er Spøgen, hvormed den listige Betragter stiller sig selv udenfor; men selv om Betragtningen brugte Rædselens Billeder for at skildre Døden, forfærdede en syg Indbildning, det er dog kun Spøg, hvis han blot tænker Døden, ikke sig selv i Døden, hvis han tænker den som Slægtens Vilkaar, men ikke som sit. Spøgen er, at hiin ubøielige Magt ligesom ikke kan faae Ram paa sit Bytte, at der er en Modsigelse, at Døden ligesom narrer sig selv. Thi Sorgen, hvis Du dermed vil sammenligne Døden, og hvis Du vil kalde den en Bueskytte som Døden er det: Sorgen rammer ikke feil, thi den rammer den Levende, og naar den har rammet ham, saa først begynder Sorgen: men naar Dødens Piil har rammet, saa er det jo forbi. Og Sygdommen, hvis Du dermed vil sammenligne Døden, og hvis Du vil kalde den en Snare, som jo Døden er den Snare, hvori Livet fanges: Sygdommen fanger virkeligt, og naar den har fanget den Karske, saa begynder Sygdommen: men naar Døden trækker Snaren sammen, saa har den jo Intet fanget, thi saa er det forbi. Men heri netop ligger Alvoren, og herved netop er Dødens Alvor forskjellig fra Livets, der saa let lader En bedrage sig selv. Thi naar Nogen gaaer sin Gang bøiet i Modgang, i Lidelser, i Sygdom, i Miskjendelse, i trange Kaar, i kummerlige Udsigter, da slutter han feil, om han ligefrem deraf slutter, at han er alvorlig; thi Alvor er ikke den ligefremme Gjengivelse, men den forædlede, det er, atter her er det det Indvortes og Tanken og Tilegnelsen og Forædlingen, der er Alvoren. Eller naar man er travl i den vidtløftige Bedrift, maaskee har mange Stridsmænd at befale over, maaskee skriver mange Bøger, maaskee er i de høie Stillinger, naar man maaskee har mange Børn, eller ofte maa være med i Livsfare, eller har den alvorlige Gjerning, at klæde Liig, saa slutter man feil, hvis man deraf uden videre slutter, at man er alvorlig, thi Alvoren er i Indtrykket, Alvoren er det indvortes Menneskes, ikke Bestillingens. Døden derimod er ikke i den Forstand noget Virkeligt, og naar man først er død, er det bag efter med at blive alvorlig, og naar man fandt en brad Død, hvilket en alvorligere Tid betragtede som den største Ulykke, hvorfor den ogsaa nævnes i den gamle Bøn, hvilket en nyere Tid anseer for den største Lykke, saa var man jo hjulpen. Livets Alvor er alvorlig, og dog er der ingen Alvor uden Bevidsthedens Forædling af det Udvortes, heri ligger Skuffelsens Mulighed; Dødens Alvor er uden Bedrag, thi det er ikke Døden der er alvorlig, men Tanken om Døden.

Naar Du derfor, m. T., vil fastholde Tanken og ikke paa anden Maade bekymre Dig om Betragtningen end ved at tænke paa Dig selv, saa skal ogsaa den umyndige Tale ved Dig blive en alvorlig Sag. At tænke sig selv død er Alvoren; at være Vidne til en Andens Død er Stemning. Det er det lette Anstrøg af Veemod, naar den Forbigaaende er en Fader, der for sidste Gang bærer sit Barn, da han bærer det til Graven; eller naar den fattige Liigvogn kjører forbi, og Du Intet veed om den Døde, kun at han var et Menneske; det er Veemod, naar Ungdom og Sundhed bliver Dødens Bytte, naar mange Aar efter Billedet af den Skjønne staaer paa det forfaldne Mindesmærke over Graven, omgivet af Ukrudt; det er Alvor i Stemning, naar Døden griber ind i de forfængelige Idrætter og griber den Daarlige, iført hendes forfængeligste Pynt, griber Daaren i hans forfængeligste Øieblik; det er Suk over Livets Spot, naar den Døde havde givet det visse Løfte og uden Skyld blev en Bedrager, da han blot havde glemt at Døden er det eneste Visse; det er Længsel efter det Evige, naar Døden tog og atter tog, og nu tog den sidste Udmærkede, Du kjendte; det er en Sjelesygdoms Feberhede eller dens kolde Brand, naar Nogen bliver saa fortrolig med Døden og med Tabet af de Nærmeste, at Livet bliver ham Aandsfortærelse; det er den rene Sorg, naar den Døde var Din; det er det udødelige Haabs Fødselssmerte, naar det var Din Elskede; det er Alvorens rystende Gjennembrud, naar det var Din eneste Veileder, og Eensomheden griber Dig: men om det var Dit Barn, og om det var Din Elskede, og om det var Din eneste Veileder, det er dog Stemning; og om Du gjerne selv vilde gaae i Døden for dem, ogsaa dette er Stemning; og om Du mener, at det er lettere, see, det er ogsaa Stemning. Alvoren er, at det er Døden Du tænker, og at Du saa tænker den som Dit Lod, og at Du saa gjør, hvad Døden jo ikke formaaer, at Du er og Døden ogsaa er. Thi Døden er Alvorens Læremester, men netop derpaa kjendes atter dens alvorlige Underviisning, at den overlader den Enkelte at søge sig op, for da at lære ham Alvor, som den kun læres ved Mennesket selv. Døden passer sin Gjerning i Livet, den løber ikke omkring som i en Frygtagtigs Forestilling og hvæsser Leen og skrækker Qvinder og Børn, som var dette Alvor. Nei, den siger: jeg er til, vil Nogen lære af mig, da komme han til mig. Kun saaledes beskæftiger Døden i Alvor, forøvrigt kun i Stemning ved Tankens Sindrighed, ved dens Dybsind, eller spøgefuldt i det oprømte Indfald, eller bøiende i den dybe Sorg, der dog i sit meest lidende Udtryk ikke er Alvor, thi Alvoren vilde netop lære at holde Maade med Sorgen og Klagen.

En Digter har fortalt om en Yngling, der den Nat, da Aaret skifter, drømte sig en Olding og som Olding i Drømmen skuede tilbage over et forspildt Liv; indtil han i Angest vaagnede Nytaarsmorgen ikke blot til et nyt Aar, men til et nyt Liv: saaledes vaagen at tænke Døden, at tænke hvad der jo er mere afgjørende end Oldingsalderen, der dog ogsaa har sin Tid, at tænke: det var forbi, at Alt var tabt med Livet for da i Livet at vinde Alt – det er Alvor. Der var en Keiser, som med Iagttagelse af alle ydre Skikke lod sig begrave. Hans Foretagende var maaskee kun en Stemning, men at være Vidne til sin egen Død, Vidne til at Kisten lukkes, Vidne til at Alt hvad der verdsligt og jordisk fylder Sandsen ophører i Døden: det er Alvor. At døe er jo ethvert Menneskes Lod og saaledes en meget ringe Kunst, men at kunne døe vel er jo den høieste Leveviisdom. Hvori ligger Forskjellen? Deri, at i det ene Tilfælde er Alvoren Dødens, i det andet den Dødeliges. Og Talen, som gjør Forskjellen, kan jo ikke henvende sig til den Døde, men til den Levende.

Saa skal da Talen handle om:

\begin{center}\emph{Dødens Afgjørelse.}\end{center}

Og deri ere vi jo enige, m. T., at en gudelig Tale aldrig bør være splidagtig, eller ueens med Andet end hvad der er ugudeligt. Naar saaledes den Fattige, den Tjenende, der sparsommeligt maa anvende den sjeldne Fridags faa Timer, gaaer ud til en Grav for at erindre en Afdød og for ogsaa at betænke sin egen Død, naar altsaa en Saadan maa hjælpe sig efter fattig Leilighed, saa Gangen derud tillige bliver en Livsnydelse, saa Opholdet derude tillige bliver en glad og velgjørende Adspredelse fra de mange Dages Arbeide, saa Tiden gaaer hen derude nu med at mindes den Afdøde, nu med at tænke paa sig selv i Alvor, nu med at glædes ved Friheden og Omgivelsen, som søgte man i en skjøn Egn Forfriskelse, som var Gangen blot til Lyst og Fødemidler derfor medbragte til forenet Glæde: da ere vi jo enige, at en Saadan i sin ædle Eenfoldighed skjønt forener Modsætninger (hvilket jo efter de Vises Ord er den høieste Vanskelighed), at hans Erindren er den Afdøde dyrebar, med Glæde annammet i Himlen, og at hans Alvor er lige saa priselig, Gud ligesaa velbehagelig, ham lige saa tjenlig som Dens, der med sjelden Evne anvendte Dag og Nat til at indøve Dødens alvorlige Tanke i sit Liv, saa han nu standsedes og atter standsedes for at forsage den forfængelige Idræt, nu tilskyndedes og atter tilskyndedes til at haste paa det Godes Vei, nu afvænnedes med at blive snaksom og travl i Livet for at lære Viisdom i Taushed, nu lærte ikke at gyse for Spøgelser og menneskelige Paafund, men for Dødens Ansvar, nu ikke at frygte dem, som ihjelslaae Legemet, men at frygte for sig selv og for at have sit Liv i Forfængelighed, i Øieblikket, i Indbildning. Vi ville prise ham, at han herligen anvendte den ham forundte Beleilighed, men hvis han derimod fra den herlige Dagens Gjerning stundom tog sig en Fridag til at forlystes ved den Tanke, at han var bedre end den Eenfoldige, der hverken havde saadan Tid eller saadan Evne, Gud velbehageligere, som gjorde Gud idel Uret, nægtede den Ene Tiden og Evnen, altsaa Lykkens Gave, og da igjen, som Mennesker stundom i Tankeløshed handle grusomt, gjorde Mangelen til en Brøde: ak, hvilken Forskjel mellem hans sjeldne Fridag og den Eenfoldiges, naar han forspilder Alt, og den Eenfoldige vinder Alt! Nei, al Sammenligning er dog Spøg, og en forfængelig Sammenligning en sørgelig Spøg. Havde end hiin Begunstigede god Tid, Alvoren og Døden vilde dog lære ham, at han ingen Tid havde at spilde, end mindre til at forspilde Alt. Skulde derimod Nogen hurtigt være bleven færdig ogsaa med Dødens Tanke som med alle andre Tanker, og fornemt maaskee bekymres for, at der i dette fattige og eensformige Liv ikke blev nok at tænke over for en saadan hurtig Tænker, da ere vi jo enige, m. T., at dette er det Eiendommelige ved enhver Gjenstand, naar den bliver det for en gudelig Betragtning, at den Eenfoldige hurtigt hjælpes til den gavnlige Forstaaelse, og at den meest Begavede glad anvender et heelt Liv, om han dog tilstaaer, at han hverken ganske har forstaaet den, eller ganske til Fuldkommenhed indøvet Tanken i sit Liv. Thi Den, der er uden Gud i Verden han bliver vel snart kjed af sig selv, og udtrykker dette fornemt ved at være kjed af hele Livet, men Den der er i Selskab med Gud, han lever jo sammen med Den, hvis Nærværelse giver selv det Ubetydeligste uendelig Betydning.

Om Dødens Afgjørelse maa der nu først siges, at den er afgjørende. Ordets Gjentagelse er det Betegnende, og Gjentagelsen selv minder om, hvor ordknap Døden er. Der er mangen anden Afgjørelse i Livet, men kun een er saaledes afgjørende som Dødens. Thi alle Livets Kræfter formaae ikke at gjøre Tiden Modstand, den river dem med sig, selv Erindringen er i det Nærværende. Og den Levende har det ikke i sin Magt at standse Tiden, at finde Hvilen udenfor Tiden i den fuldkomne Afslutning, i en Glædens Afslutning, som var der intet Imorgen, i Sorgens, som kunde den ikke blive en Draabe bittrere, i Betragtningens, som var Meningen ganske ude og ikke Betragtningen atter en Deel af Meningen, i Regnskabets, som androg Regnskabets Øieblik ikke ogsaa sit Ansvar. Døden derimod har denne Magt; den fusker ikke derpaa, som blev der dog lidt tilovers, den jager ikke efter Afgjørelsen som den Levende gjør, den gjør Alvor deraf. Naar den kommer, saa hedder det: her til, ikke et Skridt længere; saa er der sluttet af, ikke et Bogstav føies der til; saa er Meningen ude, ikke en Lyd mere skal der høres – saa er det forbi. Er det umuligt at ene alle de utallige Levendes Udsagn om Livet i eet, alle de Døde enes i eet Udsagn, i et eneste til den Levende: stat stille. Er det umuligt at ene alle de utallige Levendes Udsagn om deres Livs Stræben i eet, alle de Døde enes i eet, i eet eneste: nu er det forbi.

See, det formaaer Døden. Den er ei heller en uerfaren Yngling, der ikke har lært at bruge Leen, at Nogen skulde kunne forbløffe ham. Hav hvilken Forestilling Du vil, indbildt eller sand, om Dit Liv, om dets Vigtighed for Alle, om dets Vigtighed for Dig selv: Døden har ingen Forestilling og agter ikke paa Forestillinger. O, skulde Nogen være træt af Gjentagelsen, da vel Døden, der har seet Alt og atter og atter det Samme. Selv den i Aarhundreder sjeldne Død har den seet mange Gange; derimod har aldrig nogen Døende seet Døden skifte Farve, seet den rystet ved Synet, seet Leen vakle i hans Haand, seet Anelse om en Mineforandring i hans rolige Ansigt. Og Døden er ei heller nu bleven en gammel Mand, der svækket af Alder famler ubestemt, der ikke nøiagtigt veed hvad Klokken er, eller blev medlidende af Svaghed. O, dersom Nogen tør rose sig af at være uforandret, da vel Døden: den bliver ikke blegere og ikke ældre.

Dog skal Talen jo ikke lovprise Døden, saa lidet som den skal sysselsætte Indbildningen. At Døden kan gjøre det af, er jo vist, men Alvorens Opfordring til den Levende er at tænke det, at tænke at det er forbi, at der kommer en Tid da det er forbi. See, dette er vanskeligt; thi selv i Dødens Øieblik synes det vel den Døende, at han kunde have nogen Tid at leve i endnu, og man frygter jo endog at sige ham, at det er forbi. Og nu den Levende, saa længe han lever maaskee i Sundhed, i Ungdom, i Lykke, i Magt – betrygget altsaa, ja, vel betrygget, hvis han ikke vil lukke sig inde med Dødens Tanke, der forklarer ham, at denne Tryghed er Svig. Der er en Trøst i Livet, en falsk Smigrer, der er en Betryggelse i Livet, en hykkelsk Bedrager, den hedder: Opsættelse. Men den nævnes sjeldent ved sit rette Navn, thi selv naar man vil nævne den sniger den sig ind i Ordet, og Navnet bliver lidt formildet, og det formildede Navn er jo ogsaa en Opsættelse. Derimod er der Ingen, der saaledes kan lære at væmmes ved Smigreren og at gjennemskue Bedrageren, som Dødens alvorlige Tanke. Thi Død og Opsættelse enes ikke, de ere Dødsfjender, men den Alvorlige veed, at Døden er den Stærkeste.

Saa er det forbi. Om det var Barnet med Fordringen til et heelt Liv, om det græd for sig – nu er det forbi, ikke et Øieblik indrømmes. Om det var Ynglingen med sine skjønne Forhaabninger, om han bad for sig blot om en eneste – nu er det forbi, ikke en Hviid betales der ham paa hans Fordring til Livet. Om der manglede et Lidet i Mandens berømmelige Værk, og om dette Værk var et Verdens Underværk, og om den ganske Menneskehed vilde misforstaae det, fordi Slutningen manglede – nu er det forbi, Arbeidet ikke ganske færdigt. Om det var et eneste Ord, der havde været ham et Livs Betydning, om han vilde give et heelt Liv for at turde sige det, – nu er det forbi, Ordet blev ikke sagt.

Ved Dødens Afgjørelse er det saaledes forbi, der er Hvile; Intet, Intet forstyrrer den Døde; om hiint lille Ord, om hiint manglende Øieblik gjorde Dødens Strid urolig, nu forstyrres den Døde ikke; om hiint lille Ords Fortielse forstyrrede mange Levendes Liv, om hiint gaadefulde Værk atter og atter beskæftigede Granskeren: den Døde forstyrres ikke. Saaledes er Dødens Afgjørelse som en Nat, den Nat, der kommer, da man ikke kan arbeide; og saaledes har man jo ogsaa kaldet Døden en Nat, og formildet Forestillingen ved at kalde den en Søvn. Og det skal være formildende for den Levende, naar han søvnløs forgjeves søger Hvile paa Natteleiet, naar han undflyende sig selv forgjeves søger et Skjulested, hvor Bevidstheden ikke opdager ham, naar den Plagede træt paa Legeme og Sjel af den anstrængte Lidelse forgjeves søger en Stilling, hvori der er Lindring, der han ikke kan staae stille for Smertens Uro, og ikke gaae for Mathed, indtil han synker om og da i ny Anstrængelse forgjeves søger den hvilende Stilling, forgjeves Kølighed i denne Hede: det skal være formildende at tænke, at der dog er een Stilling, i hvilken den Anstrengede finder Hvile, det er Dødens, eet Leie, hvor han hviler stille, Dødens, een Søvn, som ikke flygter, Dødens, eet kjøligt Sted, Graven, eet Skjulested, hvor Bevidstheden staaer udenfor, Graven, hvor Erindringen selv bliver udenfor som en Luftning i Træerne, eet Sengetæppe den stille Mand ikke kaster af, hvorunder han sover roligt, Grønsværdets! Det skal være formildende, naar man i Ungdommen allerede er bleven træt, og Veemoden vil til at pusle Barnet, da at betænke, at man i Jordens Skjød ligger luunt og godt, det skal være saa formildende, at betænke denne Trøst og tænke den saaledes, at den Evige tilsidst bliver den Ulykkelige, der som en Vaagekone ikke tør sove medens alle vi Andre sove ind!

Men, m. T., dette er Stemning, og at tænke Døden saaledes er ikke Alvor. Det er Tungsindets Udflugt af Livet saaledes at længes efter Døden, og det er Oprør saaledes ikke at ville frygte den; det er Veemodets Svigagtighed ikke at ville forstaae, at der er Andet at frygte end Livet, og at der derfor maa findes en anden Viisdom, der trøster, end Dødens Søvn. Sandeligen, er det svagt at frygte Døden: da er det et opsminket Mod, der indbilder sig ikke at frygte Døden, naar det samme Menneske frygter Livet; det er Qvinde-Dorskhed, der vil til Sengs, qvindagtigt nemlig at ville sove sig til Trøst, qvindagtigt at ville sove sig fra Lidelse.

Ja vel er Døden en Søvn, og saaledes ville vi sige om Enhver, der hviler i Døden, at han sover, vi ville sige, at en stille Nat skygger over dem, og at Intet forstyrrer deres Fred. Men er der ingen Forskjel paa Liv og Død; og den Levende, der betænker sin egen Død, han betragte det anderledes. Hvis det var Dig selv, og Du den Levende, der saae det! See, Den, som sover i Døden, han rødmer ikke som Barnet i Søvne; han samler ikke ny Kraft, som Manden, der styrkes; Drømmen besøger ham ikke med sin Venlighed som den besøger Oldingen i Søvne! Naar Du i Livet seer et Tilfælde, som ligner Dødens, hvad gjør Du saa? Du raaber paa den Besvimede, fordi Du dog gyser ved denne Tilstand, det er, naar Dødens Tilstand er en Levendes; er det da trøsteligt, at Du af den Grund ikke raaber paa den Døde, fordi det ikke kan hjælpe! Men Du er jo ikke død, og vil Tungsindet styrke Dig ved Liigfaldet, vil Veemodet lade Dig besvime i en Dødens Mathed, der finder den eneste Trøst i Dødens Søvn: da raab, da kald paa Dig selv, gjør dog for Dig selv hvad Du vilde gjøre for enhver Anden, og søg ikke bedragersk Trøst i at ønske det var forbi! Hav hvilkensomhelst, indbildt eller sand Forestilling om Din Lidelses Mærkværdighed: o, dersom Nogen maatte være træt af Klageskrigets Gjentagelse, da vel Døden; selv den i Aarhundreder ved sin Lidelse sjeldne Ulykkelige, selv hans Klageskrig har Døden hørt mange Gange, men Ingen, Ingen har ymtet om, at det bevægede Døden til at komme hurtigere! Og om Dit Skrig kunde bevæge ham – er det ogsaa virkelig Din Mening, eller er det ikke snarere Modsigelsen, at han dog ikke kommer fordi Du kalder, der styrker Trodsens Selvfølelse, Modsigelsen der hjælper den Frygtagtige til at lege den modige Leg med den Forfærdelige – om altsaa Dit Skrig og Din Længsel bevægede ham, mon Du dog ikke bedrog Dig selv, om vi end et Øieblik ville glemme Ansvaret som altid bliver? Hvad var det der lindrede, mon det var det, at det var forbi, eller ikke Forestillingen derom, som denne endnu var i Tungsindets og Veemodets og saaledes i den jo Levendes Magt – en Adspredelse, et Legetøi! See, Den, som sover i Døden, han rører sig ikke, og hvis end Kistens Klædning ikke sluttede saa stramt om ham – han rører sig dog ikke; han vorder til Støv. Og Tanken om at det er forbi, der i Forestillingens indbildte Forskud tungsindigt vederqvægede i trodsig Afmagt, eller leflende lindrede i Veemod, den er der jo ikke hos ham. Han har altsaa ingen Glæde af, at det er forbi: hvorfor ønskede han det da saa meget? Hvilken Modsigelse! Siig saa, at det er saa trøsteligt at forraadne i Jorden. Men veed Du Andet om Døden, da veed Du ogsaa, at frygte Andet end Livet.

Alvoren forstaaer vel det Samme om Døden, men forstaaer det anderledes. Den forstaaer, at det er forbi. Om dette lader sig, formildet i Stemning, udtrykke saaledes, at Døden er en Nat, en Søvn, beskæftiger den mindre. Alvoren spilder ikke megen Tid paa at gjette Gaader, den sidder ikke hensjunken i Betragtningen, omskriver ikke Udtrykkene, betænker ikke Billedsprogets Sindrighed, afhandler ikke, men handler. Er det vist, at Døden er til, som det er; er det vist, at det med dens Afgjørelse er forbi; er det vist at Døden selv aldrig indlader sig paa at give nogen Forklaring: nu vel, da gjelder det at forstaae sig selv, og Alvorens Forstaaelse er, at er Døden Natten, saa er Livet Dagen, kan der ikke arbeides om Natten, saa kan der arbeides om Dagen; og Alvorens korte men tilskyndende Raab, ligesom Dødens korte, er: endnu idag. Thi Døden i Alvor giver Livskraft som intet Andet, den gjør aarvaagen som intet Andet. Det sandselige Menneske virker Døden paa, til at sige: lader os æde og drikke, thi imorgen skulle vi døe; men dette er Sandselighedens feige Livslyst, hiin foragtelige Tingenes Orden, hvor man lever for at spise og drikke, ikke spiser og drikker for at leve. Det dybere Menneske virker Dødens Forestilling maaskee paa til Afmagt, saa han segner slappet i Stemning; men den Alvorlige giver Dødens Tanke den rette Fart i Livet og det rette Maal, hvorefter han retter Farten. Og ingen Buestræng kan strammes saaledes og formaaer at give Pilen en saadan Fart, som Dødens Tanke formaaer at fremskynde den Levende, naar Alvoren spænder den. Da griber Alvoren det Nærværende endnu idag, forsmaaer ingen Opgave som for ringe, vrager ingen Tid som for kort, arbeider af yderste Evne, om den dog er villig til at smile over sig selv, hvis denne Anstrængelse skulde være Fortjeneste for Gud, og villig til i Afmagten at forstaae, at et Menneske er slet Intet, og at Den, der arbeider af yderste Evne, kun faaer ret Leilighed til at forundre sig over Gud. Tiden er jo dog ogsaa et Gode. Formaaede et Menneske at frembringe Dyrtid i den udvortes Verden, ja da vilde han faae travlt; thi Kjøbmanden siger jo rigtigt, at vel har Varen sin Priis, men Prisen afhænger dog saa meget af de fordeelagtige Tidsforhold – og naar der er Dyrtid, saa tjener Kjøbmanden. I den udvortes Verden formaaer et Menneske det nu maaskee ikke, men i Aandens Verden formaaer Enhver det. Døden frembringer jo selv Dyrtid paa Tid i Forhold til den Døende, hvo har ikke hørt, hvorledes en Dag stundom en Time blev skruet op i Priis, naar den Døende tingede med Døden; hvo har ikke hørt, hvorledes en Dag stundom en Time fik uendeligt Værd, fordi Døden gjorde Tiden dyr! Dette formaaer Døden, men den Alvorlige formaaer ved Dødens Tanke at gjøre Dyrtid, saa Aaret og Dagen faaer uendeligt Værd – og naar der er dyr Tid, saa tjener Kjøbmanden ved at bruge Tiden. Men naar saa den borgerlige Sikkerhed er foruroliget, saa dynger Kjøbmanden ikke det Bundne ligegyldigen op, men han vaager over sin Skat, at Tyvehaand ikke skal bryde ind og tage den fra ham: ak, Døden er jo ogsaa som en Tyv om Natten.

Ikke sandt, m. T., saaledes har Du selv erfaret det? Og naar Dødens Tanke besøgte Dig, men gjorde Dig uvirksom, naar den listede sig ind og bedaarede Livskraften i sværmersk Drøm; naar Dødens Mismod vilde gjøre Dig Livet til Forfængelighed; naar hiin Forfører Veemodet sneeg sig omkring Dig; naar Forestillingen, som var det forbi, vilde bedøve Dig i Tungsinds Søvn; naar Du hensank i Aandsfraværelsens Syssel med Dødens Sindbillede: da gav Du ikke Døden Skyld derfor, thi alt Dette var jo ikke Døden. Men Du sagde til Dig selv: min Sjel er i Stemning, og vedbliver det saaledes, da er deri et Fjendskab mod mig, som kan faae Overmagten. Da flyede Du ikke Døden, som var dette det Helbredende. Langtfra. Du sagde: jeg vil fremkalde Dødens alvorlige Tanke. Og den hjalp Dig. Thi Dødens Alvor har hjulpet til at gjøre en sidste Time uendeligt betydningsfuld, dens alvorlige Tanke hjulpet til at gjøre et langt Liv betydningsfuldt som i Dyrtid, varagtigt som efterstræbt af Tyvehaand.

Saa lad da Døden beholde sin Magt, »at det er forbi«, men Livet ogsaa beholde sin Ret til at arbeide medens det er Dag; og lad den Alvorlige søge Dødens Tanke som Bistand dertil. Den Vankelmodige er blot et Vidne til den bestandige Grændsestrid mellem Liv og Død, hans Liv kun Tvivlens Angivelse af Forholdet, hans Livs Udgang en Skuffelse, men den Alvorlige har sluttet Venskab med de Stridende, og hans Liv har i Dødens alvorlige Tanke den meest trofaste Forenede. Er der da end den Lighed for alle Døde, at det nu er forbi, der er dog een Forskjel, m. T., den raaber til Himlen denne Forskjel, Forskjellen paa, hvilket Liv det var, som nu med Døden er forbi. Saaledes er det altsaa ikke forbi, og trods al Dødens Forfærdelse, nei, understøttet af Dødens alvorlige Tanke siger den Alvorlige: det er ikke forbi. Men frister denne lyse Udsigt, skimter han atter blot efter den i Betragtningens Dæmring, fjerner den ham fra Opgaven, bliver Tiden ikke Dyrtid, er Besiddelsen ham tryg; da er han atter ikke alvorlig. Siger Døden: maaskee endnu idag; da siger Alvoren: det være nu maaskee idag eller ikke, jeg siger: endnu idag.

Om Dødens Afgjørelse maa dernæst siges, at den er ubestemmelig. Hermed er Intet sagt, men saaledes maa det ogsaa være, naar Talen er om en Gaade. Vel gjør Døden nemlig Alle lige, men dersom denne Lighed er i Intet, i Tilintetgjorthed, saa er jo Ligheden selv ubestemmelig. Skal der tales videre om denne Lighed, da kan det kun skee ved at nævne Livets Forskjellighed og benegte denne om Dødens Lighed. Her, i Graven, er Barnet og Den som omskabte en Verden, lige uvirksomme; her er den Rige ligesaa fattig som den Fattige, Armoden trygler ikke, den Rige har Intet at give af, den Nøisomste og den Umætteligste behøve lige Lidet; her høres ikke Herskerens Stemme, og ikke den Undertryktes Skrig; her blev den Overmodige og den Krænkede lige afmægtige; her ligge de Grav om Grav og taale hinanden, de, hvem Fjendskabet adskilte med en Verden; her ligger den Deilige og her den Elendige, men Deiligheden adskiller dem ikke; her ligge de Begge, Den der speidede efter Døden som efter en skjult Skat, og Den som havde glemt at Døden var til, men Forskjellen er ikke til at opdage.

Saaledes er Dødens Afgjørelse ved sin Lighed som det tomme Rum og som en Taushed, hvori Intet lyder, eller formildet som en Taushed, der ikke forstyrres. Og i dette tause Rige hersker Døden. Skjøndt een mod alle de Levende, er den dog mægtig til at underlægge sig dem og til at paabyde Taushed. Hav hvilken Forestilling Du vil om Dit Liv, ja selv om dets Betydning for det Evige, Du taler Dig ikke fra Døden, Du gjør ikke Overgangen til det Evige i Talens Løb og med eet Aandedrag: de have Alle maattet tie. Og forenede end Slægt med Slægt sig larmende til fælles Gjerning, og glemte den Enkelte sig selv og fandt sig saa betrygget i Mængdens Skjul: see, Døden tager Hver for sig – og han bliver taus. Var end den Levendes Forskjel hvilken Du vil tænke, Døden gjør ham lige med Den, der var ukjendelig ved sin Forskjellighed. Thi Livets Speil gjengiver vel stundom den Forfængelige med smigrende Troskab hans Forskjellighed, men Dødens Speil smigrer ikke, dets Troskab viser Alle eens, de ligne Alle hinanden, naar Døden med sit Speil har prøvet, at den Døde tier.

Saaledes er Dødens Afgjørelse ubestemmelig ved Ligheden, thi Ligheden er i Tilintetgjorthed. Og det skal være formildende at tænke derpaa for den Levende. Naar Aanden, træt af Forskjelligheden, som bliver ved og bliver ved og aldrig bliver færdig, stolt trækker sig tilbage i sig selv og samler Vrede i Afmagtens Trods, at den ikke formaaer at standse Forskjellighedens Livskraft: da skal det være formildende, at betænke, at Døden har denne Magt, da skal denne Forestilling puste hiin Tilintetgjørelsens Begeistring op til en Gløden, hvori der skal være forhøiet Liv. – Naar den Elendige sukker i sin Afkrog, fordi Livet for urettede ham som en Stedmoder, naar han vanheldet end ikke tør vise sig, fordi selv det bedste Menneske uvilkaarligt smiler af hans qvalfulde, ak og dog latterlige Jammer, naar han saaledes i afsondret Udelukke ikke elsker, fordi Ingen finder det Lige hos ham, hvad han selv forgjæves søger hos de Andre: da skal det være lindrende som en Køling af Snee for den skjulte Harmes Brand, at betænke, at Døden gjør Alle lige. – Naar den Krænkede vrider sig under den Mægtiges Uret, og Hadet i Afmagt fortvivler om Hævn: da skal det være den velkomne Trøst, der næsten kalder Livslysten tilbage, at Døden gjør dem Alle lige. – Naar den ved Ønsker Forkjelede sidder uvirksom og lefler med Ønskets store Forestillinger om sig selv, men kun seer Andre stræbe og naae det Store, naar saa Utaalmodighedens Lidenskab besværer den Forkjeledes Aandedræt: da skal det være lindrende, det skal give Luft, at betænke, at Døden slaaer en Streg over det Hele og gjør Alle lige. – Naar den Tabende vel forstod, at Striden nu var forbi, og han den Svagere, men tillige forstaaer, at det dog ikke er forbi, at hans Nederlag gav Seierherren Lykkens Fart, at hans Liden i Nederlagets Efterfølge dagligt, men fjernere og fjernere er Efterretningen om den Andens Stigen i det Fjerne: da skal det være formildende, at betænke, at Døden henter ham ind og gjør Adskillelsen til et Intet. – Naar Sygdom bliver den daglige Gjest, og Tiden gaaer hen, Glædens Tid, naar selv de Nærmeste blive trætte af den Lidende og mangt et utaalmodigt Ord saarer, naar den Lidende selv føler, at blot hans Nærværelse er forstyrrende for de Glade, saa han maa sidde fjern fra Dandsen: da skal det være lindrende at betænke, at Døden indbyder dog ogsaa ham til Dands og at i den Dands blive Alle lige.

Men, m. T., dette er Stemning; og egentligen er det Feighed, der ved et Falskmaal, iført digterisk Skikkelse, vil tykkes sig bedre, om den dog i Væsen er ligesaa ussel. Thi om den Eenfoldige maaskee ikke er istand til at fatte den Art Stemning, er denne Fornemhed i og for sig et afgjørende Værd, er den ikke kun afgjørende til at gjøre den mere tilregnelig? Det er Tungsinds feige Lyst, at ville svimle sig ind i det Tomme, og derved søge i denne Svimmel den sidste Adspredelse; det er Misundelse i Oprør mod Gud at ville tage Skade paa sin Sjel, saaret af Forskjelligheden; det er Selvangivelse, afmægtigt at ville hade, den forraader, at man blot mangler Magten, da man jo gjør den forfærdelige Misbrug af Afmagten; det er en foragtelig Gjenvei til ubeføiet Klage over Livet, at man blot vil ønske, og da klage fordi man ikke blev hvad man ønskede, og aldrig blive duelig til Andet end til at ønske, og endeligen elendig nok til at ønske Alt borte; det er selvplagersk Udholdenhed af den Overvundne, intet Høiere at ville forstaae end Striden mellem Dig og Mig og begges Undergang; det er en endnu forfærdeligere Sygdom ikke at ville fatte, hvilken Læge den Syge har Behov. Sandeligen, er det feigt og vellystigt Kjelenskab, end ikke i Tanken at turde opgive den begunstigende Forskjellighed og at have sit Liv fortabt i den, da er det et opsminket Mod, der vil forsøge sig paa Forestillingen om Dødens Lighed, naar dog det samme Menneske sukker eller gisper under Livets Forskjellighed.

Og hvis det virkelig var Nogens Mening – skulde det ikke være Modsigelsen, at han endnu levede, der gav det formastelige Vovestykke sin Tillokkelse – at ville trøstes saaledes ved Dødens Lighed, mon hans Forestilling om Døden holdt Stik i Døden, det er, naar Tankesysselen ikke mere forlystede Lidenskaben? Den Døde har jo glemt Forskjelligheden; og om han end gjennem et heelt Liv foresatte sig at ville erindre den, for at have Glæde af at see den fratagen en Anden i Døden, i Døden er jo denne Tanke ikke hos ham, selv om vi i et Øieblik ville glemme Ansvaret, som venter. Dette er Løgnen og Bedraget i den formastelige Trods, der vil sammensværge sig med Døden mod Livet. Det glemmes, at Døden er den Stærkeste, det glemmes, at den er uden Forkjerlighed, at den ikke slutter Forbund med Nogen, saa han i Døden fik Frispas og Raaderum til Tilintetgjørelsens Lyst. Kun naar den Levendes Forestilling eventyrligt vanker om i de Dødes tause Rige, gøglende selv er Døden, og forsvinder for sig selv i Døden; kun naar den Levendes Forestilling leger Døden efter, stævner den Misundte ind, afklæder ham al hans Herlighed og fryder sig ved hans Afmagt; kun naar Forestillingen gaaer ud til Gravene, sætter formasteligt Spaden i Jorden, krænker de Dødes Fred ved den trodsige Lyst, at den ene Afsjeledes Levninger see ud aldeles som den Andens – kun da lindrer det.

Men alt Sligt er ikke Alvor; og er dets Væsen end nok saa mørkt, og er Forlystelsen end nok saa skummel, derfor er det ikke Alvor. Thi Alvor skuler ikke, men er forliget med Livet og veed at frygte Døden.

Alvoren forstaaer da det Samme om Døden, men forstaaer det anderledes. Den forstaaer, at Døden gjør Alle lige; og det har den allerede forstaaet, fordi Alvoren har lært den at søge Ligheden for Gud, i hvilken Alle kunne være lige. Og i denne Stræben opdager den Alvorlige en Forskjellighed, sin egen nemlig fra det Maal, som er sat ham, og opdager, at fjernest fra dette Maal vilde en Tilstand være som den Lighed, hvilken er Dødens. Men hver Gang den jordiske Forskjellighed vil friste, vil sinke, da træder den alvorlige Tanke om Dødens Lighed imellem og tilskynder atter. Som ingen ond Aand tør nævne det hellige Navn, saaledes gyser enhver god Aand for det tomme Rum, for Tilintetgjorthedens Lighed, og denne Gysen, der er frembringende i Naturens Liv, er fremskyndende i Aandens. O, hvor ofte lærte ikke Tilintetgjørelsens Lighed, naar Døden kom til et Menneske, ham at ønske den tungeste Forskjellighed tilbage, at finde Vilkaaret ønskeligt, nu da Dødens Vilkaar var det eneste! Og saaledes har Dødens alvorlige Tanke lært den Levende at gjennemtrænge den tungeste Forskjellighed med Ligheden for Gud. Og ingen Sammenligning har den fremskyndende Magt, og giver den Hastende saa sikkert den sande Retning, som naar den Levende sammenligner sig selv med Dødens Lighed. Og er af al Sammenligning den den forfængeligste, naar et Menneske forsmaaer enhver anden Sammenligning for at sammenligne sig med sig selv i Selvtilfredshed, ja stod maaskee ingen forfængelig Qvinde saa forfængelig omgiven af Beundringen, som da hun stod eensom for sit Speil: o, saa er der ingen Sammenligning saa alvorlig, som Dens, der eensom sammenlignede sig med Dødens Lighed. Eensom; thi det er jo ogsaa det som Dødens Lighed gjør ham, naar Graven lukkes, naar Porten lukkes for Haven, naar Natten falder paa, og han ligger eensom, fjern fra al Deeltagelse, ukjendelig, i den Skikkelse som kun kan vække Gysen, eensom derude, hvor de Dødes Mængde ikke danner noget Selskab. See, Døden har formaaet at omstyrte Throner og Fyrstendømmer, men Dødens alvorlige Tanke har gjort det lige saa Store, har hjulpet den Alvorlige til at underlægge den meest begunstigede Forskjellighed under den ydmyge Lighed for Gud, og hjulpet til at opløfte sig over den tungeste Forskjellighed i den ydmyge Lighed for Gud.

Ikke sandt, m. T., saaledes har Du selv erfaret det. Og naar Din Sjel foer vild i Begunstigelse, og naar Du neppe kunde kjende Dig selv igjen for Herligheden, da gjorde den alvorlige Tanke om Dødens Lighed Dig paa en anden Maade ukjendelig, og Du lærte at kjende Dig selv, og at ville være kjendt for Gud. Eller dersom Din Sjel sukkede i Lidelsers, i Gjenvordigheds, i Krænkelses, i Tungsinds svære Indskrænkning, ak, og det syntes Dig at Indskrænkningen vilde blive livsvarig, naar da Fristeren ogsaa kom til dit Huus, Du veed, han, Fristeren, som man har i sit eget Indre, og som bedragersk hilser fra Andre, og naar han saa først foregøglede Dig de Andres Lykke, indtil Du blev mismodig, og han da vilde give Dig Opreisning: da hengav Du Dig ikke i Stemningen. Du sagde, den er et Oprør mod Gud, et Fjendskab mod mig selv; og da sagde Du, jeg vil kalde Dødens alvorlige Tanke frem. Og den hjalp Dig til at overvinde Forskjelligheden, til at finde Ligheden for Gud, til at ville udtrykke Ligheden. Thi Dødens Lighed er forfærdelig derved, at Intet kan modstaae den (hvor trøstesløst!), men Guds-Ligheden er salig derved, at Intet kan forhindre den, hvis Mennesket ikke selv vil. Og var da saa Livets Forskjellighed saa stor! Thi tag den Glade, lad ham glæde sig ved sin Lykke, naar Du, den Ulykkelige, var glad igjen ved hans Lykke, da vare I jo begge Glade? Tag den Udmærkede, lad ham fryde sig ved sit Fortrin, naar Du, den Krænkede, havde glemt Fornærmelsen, og nu saae det Fortrinlige hos ham, var Forskjellen vel stor? Tag den Unge, lad ham ile frem med Haabets Tilforladelighed, naar Du, skjøndt skuffet af Livet, maaskee skjult endog understøttede ham, var Forskjellen da saa stor? O, Lykke og Ære og Rigdom og Skjønhed og Magt, det er jo hvad der udgjør Forskjelligheden, men hvis Forskjellen kun er den, at den Enes Lykke og Ære og Rigdom og Skjønhed og Magt er en Frilands-Plante, den Andens en Gravblomst der dyrkes i Selvfornægtelsens helligede Jord: er Forskjellen da saa stor, de ere jo begge lykkelige, og ærede, og rige, og skjønne og mægtige. Ak nei, da behøver man ingen Opreisning, mindst en saadan, der falsk fortier, at man selv bliver til Intet. Hvor tung end Forskjelligheden var, den alvorlige Tanke om Dødens Lighed hjalp dog som den strænge Opdragelse til at forsage verdslig Sammenligning, til at forstaae Tilintetgjørelsen som det endnu Forfærdeligere, og til at ville søge Ligheden for Gud.

Dødens Lighed fik ikke Lov til at fortrylle Dig med sin Trolddom; der er jo ei heller Tid dertil. Thi som Dødens Afgjørelse er ubestemmelig ved Ligheden, saa er den lige saa ubestemmelig ved Uligheden. Hvo har saaledes ikke ofte hørt Talen om, at Døden ingen Forskjel gjør, at den ikke kjender Stand og ikke Alder; hvo har ikke selv ofte overveiet det, at naar han nævnede det forskjelligste en Levendes Vilkaar, og han nu vilde tænke Døden i Forhold dertil, at dens Bestemmelse da var, at den ligesaa godt kunde søge sig Bytte her som der, lige saa godt, fordi der intet Hensyn tages, medens al Forskjellighed netop ligger i at tage Hensyn. Saaledes er den ubestemmelig ved sin Ulighed. Den kommer næsten Livet i Forkjøbet, og Barnet fødes dødt, den lader Oldingen vente fra Aar til Aar; naar man siger Fred og Tryghed, da staaer den over Een, og i Livsfaren søges den stundom forgjeves, medens den finder Den, som skjuler sig i en Afkrog; naar Laderne ere fulde og der er Forraad for et langt Liv, saa kommer Døden og kræver den Riges Sjel, naar der er Mangel, bliver den borte; naar den Hungrige bekymret sørger for, hvad han skal faae at spise imorgen, saa kommer Døden og tager Næringssorgerne fra ham, naar Vellystningen overmættet bekymres for, hvad han skal faae at spise imorgen, saa kommer Døden dømmende og gjør Bekymringen overflødig.

Saaledes er Døden ubestemmelig: det eneste Visse, og det Eneste, hvorom det Intet er vist. Denne Forestilling lokker Tanken ud i det Ubestemmeliges Vexel for at ville forsøge sig i denne Gysen som i en Leeg, for at ville gjette denne forunderlige Gaade, for at ville hengive sig i det Pludseliges uforklarlige Forsvinden og uforklarlige Frembryden. Det skal være formildende at tænke over dette Træf, dette Lige og Ulige, denne anede Lov i det Lovløse, som er og som er ikke, er i Forhold til alt Levende, og som er ubestemmelig i ethvert sit Forhold. Naar Sjelen bliver træt af Tvang og Bundethed, af det Bestemmelige og af den bestemmelige Opgaves knappe daglige Maal, og af Bevidstheden om, at der forsømmes mere og mere; naar Villiens Kraft har tjent ud, og den Udmarvede bliver ligesom Trøske; naar Nysgjerrigheden træt af Livet søger en mangfoldigere Opgave for Nysgjerrigheden: da skal det være underholdende at betænke Dødens Ubestemmelighed, og formildende saaledes at blive fortrolig med denne Tanke. Nu forundrer man sig over eet Dødsfald, nu over et andet, nu taler man sig ør ved i almindelige Udtryk at tale om, hvad der unddrager sig den almindelige Bestemmelse, nu er man i een Stemning, nu i en anden, nu veemodig, nu uforfærdet, nu spottende, nu knyttende Døden til det lykkeligste Øieblik som den største Lykke, nu som den største Ulykke, nu ønskende en brad Død, nu en langsom, nu trættes man i Samtale om, hvilken Død der dog er det ønskeligste, nu bliver man kjed af hele Overveielsen, glemmer Døden, indtil Betragtningens Hjul atter sættes i Bevægelse og ryster Betragtningens Enkeltheder sammen i nye Forbindelser til ny Forundring – ak ja, indtil Tanken om den egne Død fordunster i en Taage for Øiet, og Paamindelsen om den egne Død bliver en ubestemt Susen for Øret. Denne er Fortrolighedens Formildelse i den Afsløvedes Betragtning, at det nu engang er saaledes, i den opløftende upersonlige Glemsomhed, der glemmer sig selv over det Hele, eller rettere sig selv i Tankeløshed, hvorved den egne Død bliver et snurrigt Tilfælde med i disse mangfoldige uberegnelige Tilfælde, og Udtjentheden en Forberedelse, der gjør den egne Døds Overgang mild.

Men selv om et saadant Liv ved at betænke Dødens Forunderlighed gjennemgik alle mulige Stemninger, er Betragtningen derfor Alvor? Ender Stemningers Vidtløftighed altid i Alvor, skulde Alvorens Begyndelse ikke snarere være at forhindre hiin Vidtløftighed, i hvilken den Betragtende forsømmer Livet og bliver som Den, der forfalder til Spil, naar han grunder og grunder og drømmer om Tal i Natten, istedenfor at arbeide om Dagen. Den, der betragter Døden saaledes, er i en bedøvet Tilstand med Hensyn til sit aandelige Liv; han svækker sin Bevidsthed, saa den ikke kan udholde det alvorlige Indtryk af det Uforklarlige, saa han ikke i Alvor kan bøie sig under Indtrykket, men da ogsaa betvinge den Gaadefulde.

Ja vist er Døden en forunderlig Gaade, men kun Alvoren kan bestemme den. Hvoraf kommer vel hiin Tankeløshedens Forvirring uden deraf, at den Enkeltes Tanke vover sig betragtende ud i Livet, vil overskue hele Tilværelsen, hiint Spil af Kræfter som kun Gud i Himlene kan roligt betragte, fordi han i sit Forsyn behersker det med den vise og allestedsnærværende Hensigt, men som svækker et Menneskes Aand, og gjør ham sindssvag, volder ham utidig Sorg, og styrker med beklagelig Trøst. Utidig Sorg nemlig i Stemning, fordi han bekymrer sig om saa Meget, beklagelig Trøst nemlig i afspændt Sløvhed, naar hans Betragtning har saa mange Indgange og Udgange, at den tilsidst bliver en Vildelse. Og naar saa Døden kommer, da bedrager den dog Betragteren, fordi al hans Betragtning ikke kom Forklaringen et eneste Skridt nærmere, men kun bedrog ham for Livet.

Alvoren forstaaer da det Samme om Døden, at den er ubestemmelig ved Uligheden, at ingen Alder, og ingen Omstændighed, og intet Livsforhold sikkrer mod den, men derpaa forstaaer den Alvorlige det anderledes og forstaaer sig selv. See, Øxen ligger allerede ved Roden af Træet, hvert Træ, som ikke bærer god Frugt, skal omhugges – nei, ethvert Træ skal omhugges, ogsaa det, som bærer god Frugt. Det Visse er, at Øxen ligger ved Roden af Træet; om Du end ikke bemærker, at Døden gaaer over Din Grav og at Øxen bevæger sig, Uvisheden er dog i ethvert Øieblik, det Uvisse naar Hugget falder – og Træet. Men naar det er faldet, da er det afgjort, om Træet bar god Frugt eller raadden Frugt.

Den Alvorlige betragter sig selv; er han ung da lærer Tanken om Døden ham, at det er et ungt Menneske, som her bliver dens Bytte, hvis den kommer idag, men han fjanter ikke i almindelig Tale om Ungdom som Dødens Bytte. Den Alvorlige betragter sig selv, han veed altsaa, hvorledes Den er beskaffen, som her vilde blive Dødens Bytte, hvis den kom idag; han betænker sin egen Gjerning og veed altsaa, hvilken Gjerning det var, som her afbrødes, hvis Døden kom idag. Saa hører Legen op, saa er Gaaden gjettet. Den almindelige Betragtning af Døden forvirrer kun Tanken, ligesom det at ville erfare i Almindelighed. Dødens Vished er Alvoren, dens Uvished er Underviisningen, Indøvelsen af Alvoren; den Alvorlige er Den, som ved Uvisheden opdrages til Alvor i Kraft af Visheden. Hvorledes lærer vel et Menneske Alvor? Mon derved, at en Alvorlig foresiger ham Noget, at han kunde lære dette? Ingenlunde. Har Du ikke selv saaledes lært af en alvorlig Mand, saa tænk Dig, hvorledes det gaaer til. See, den Lærende bekymrer sig (thi uden Bekymring ingen Lærende) om en eller anden Gjenstand med hele sin Sjel; og saaledes er jo Dødens Vished en Bekymringens Gjenstand. Nu henvender den Bekymrede sig til Alvorens Lærer; og saaledes er jo Døden, ikke en Skræmsel uden for Indbildningen. Den Lærende vil nu Dette eller Hiint, han vil gjøre det saaledes under disse Forudsætninger: »og ikke sandt, saa lykkes det?« Men den Alvorlige svarer slet Intet, og endeligen siger han, dog uden at spotte, med Alvorens Rolighed: »ja, det er muligt!« Den Lærende bliver allerede lidt utaalmodig; han udkaster en ny Plan, forandrer Forudsætningerne, og slutter sin Tale paa en endnu mere indtrængende Maade. Men den Alvorlige tier, seer roligt paa ham og siger endeligen: »ja, det er muligt!« Nu bliver den Lærende lidenskabelig, han griber til Bønner, eller hvis han er saaledes udrustet til snilde Tankevendinger, ja han fornærmer maaskee endog den Alvorlige, og bliver selv ganske forvirret, og Alt synes Forvirring om ham; men da han med disse Vaaben og i denne Tilstand stormer ind paa den Alvorlige, maa han udholde hans uforandrede rolige Blik og finde sig i hans Taushed, thi den Alvorlige seer blot paa ham og siger endeligen: »ja det er muligt.« Saaledes med Døden. Visheden er den Uforanderlige, og Uvisheden er det korte Ord: det er muligt; og enhver Betingelse, der vil gjøre Dødens Vished til en betinget Vished for den Ønskende, enhver Overeenskomst, der vil gjøre Dødens Vished til en betinget Vished for den Besluttende, enhver Aftale, der vil betinge Dødens Vished ved Tid og Time for den Handlende, enhver Betingelse, enhver Overeenskomst, enhver Aftale strander paa dette Ord; og al Lidenskabelighed og al Snildhed og al Trods gjøres afmægtig ved dette Ord, indtil den Lærende gaaer i sig selv. Men netop deri ligger Alvoren, og netop dertil var det, at Visheden og Uvisheden vilde hjælpe den Lærende. Faaer Visheden Lov at staae hen for hvad den nu kan være, som en almindelig Overskrift over Livet, ikke, som det skeer ved Uvishedens Hjælp, som Anvendelsens Paategning paa det Enkelte og Daglige, da læres Alvoren ikke. Uvisheden træder til og peger stadigen som Læreren paa Lærens Gjenstand, og siger til den Lærende: læg vel Mærke til Visheden: da bliver Alvoren til. Og ingen Lærer formaaer saaledes at lære Discipelen at lægge Mærke til hvad der siges, som Dødens Uvished, naar den peger paa Dødens Vished; og ingen Lærer formaaer saaledes at holde Discipelens Tanker samlede paa Underviisningens ene Gjenstand, som Tanken om Dødens Uvished, naar den indøver Tanken om Dødens Vished.

Dødens Vished bestemmer den Lærende engang for alle i Alvor, men Dødens Uvished er det daglige, eller dog det hyppige, eller dog det fornødne Tilsyn, der vaager over Alvoren: først dette er Alvor. Og intet Opsyn er saa omhyggeligt, ikke Faderens med Barnet, ikke Lærerens med Lærlingen, ja ikke Fangevogterens med Fangen; og intet Tilsyn saa forædlende som Dødens Uvished, naar den prøver Tidens Anvendelse og Gjerningens Beskaffenhed, den Besluttendes eller den Handlendes, Ynglingens eller Oldingens, Mandens eller Qvindens. Thi med Hensyn til vel anvendt Tid er det i Forhold til Dødens Afbrydelse ikke væsentligt, om Tiden var lang eller kort; og med Hensyn til den væsentlige Gjerning er det i Forhold til Dødens Afbrydelse ikke væsentligt, om den blev færdig eller kun begyndt. Med Hensyn til det Tilfældige er Tidens Længde bestemmende, som for at nævne Lykken: Enden afgjør først, om man har været lykkelig. I Forhold til den tilfældige Gjerning, som er i det Udvortes, er det væsentligt, at Værket bliver færdigt. Men den væsentlige Gjerning bestemmes ikke væsentligen ved Tiden og det Udvortes, forsaavidt Døden er Afbrydelsen. Saa bliver Alvoren den at leve hver Dag, som var det den sidste og tillige den første i et langt Liv; og at vælge Gjerningen som ikke er afhængig af, om der undes en Menneskealder til at fuldkomme den vel, eller kun en kort Tid til at have begyndt den vel.

Endeligen maa der siges om Dødens Afgjørelse, at den er uforklarlig. Om nemlig Menneskene finde en Forklaring: Døden selv forklarer Intet. Thi hvis Du kunde faae Øie paa ham, den blege, glædeløse Høstkarl, hvor han stod ledig, støttende sig ved Leen, og Du da vilde gaae hen til ham, det være nu, at Du meente, Din Kjede ved Livet maatte indynde Dig hos ham, eller at Din brændende Længsel efter det Evige skulde røre ham, dersom Du lagde Din Haand paa hans Skuldre og sagde: forklar Dig, blot et Ord – troer Du, han svarede? Jeg tænker, han mærkede det end ikke, at Du lagde Haanden paa hans Skuldre og talte til ham. Eller dersom Døden kom, ak, saa beleilig, ak, som den største Velgjører, som en Frelser, dersom den kom og frelste et Menneske fra at paadrage den Skyld, som ikke bliver angret i Livet, fordi Skylden gjør Ende paa Livet, hvis nu hiin Ulykkelige vilde takke Døden, at den bragte ham det Søgte og forhindrede ham i at blive skyldig, troer Du den forstod ham? Jeg tænker, den hørte end ikke et Ord af hvad han sagde; thi den forklarer Intet. Om den kommer som den største Velgjerning eller som den største Ulykke, om den hilses med Jubel eller med fortvivlet Modstand, det veed Døden Intet af, thi den er uforklarlig. Den er Overgangen; om Forholdet veed den Intet, slet Intet.

See, denne Uforklarlighed trænger jo vel til en Forklaring. Men deri ligger netop Alvoren, at Forklaringen ikke forklarer Døden, men aabenbarer, hvorledes den Forklarende er i sit inderste Væsen. O, alvorlige Paamindelse om Langsomhed i at tale! Maa man end smile, naar man seer Tankeløsheden sætte Haanden understøttende under det grublende Hoved, der skal udgrunde Forklaringen, maa man end atter smile, naar saa denne Tænker rykker ud med Forklaringen; eller naar, som var det paa et almindeligt Opbud, selv de letfærdigste Tanker i Forbigaaende ere tilrede med et Indfald, en Bemærkning som Forklaring, benyttende den sjeldne Beleilighed, da jo Døden er Alle en uforklarlig Gaade: ak, den dømmende Alvor over slig Færd er, at den Forklarende angiver sig selv, forraader, hvor tankeløst, hvor daarligt hans Liv er. Derfor er Tilbageholdenhed med Forklaringen allerede et Tegn paa nogen Alvor, der dog forstaaer, at Døden, netop fordi den er Intet, er ikke noget Saadant, som en forunderlig Indskrift, hvilken enhver Forbigaaende skal søge at læse, eller som en Mærkværdighed, hvilken Enhver maa have seet og have en Mening om. Det Afgjørende ved Forklaringen, Det, som forhindrer, at Dødens Intet ikke gjør Forklaringen til Intet, er at den faaer tilbagevirkende Kraft og derved Virkelighed i den Levendes Liv, saa Døden bliver ham en Lærer, og ikke forrædersk hjælper ham til en Selvangivelse, der giver Forklareren an som en Daare.

Som det Uforklarlige kan jo Døden synes at være Alt og slet Intet, og Forklaringen synes at være, at udsige dette i eet. En saadan Forklaring angiver et Liv, der nøiet med det Nærværende værger sig mod Dødens Indflydelse ved en Stemning, som holder denne i Uafgjorthedens Ligevægt. Døden faaer ikke Magt til at forstyrre et saadant Liv, faaer derimod Indflydelse, men ikke tilbagevirkende Kraft til at omdanne et saadant Liv. Forklaringen vexler ikke i forskjellige Stemninger, men Døden føres hvert Øieblik udenfor Livet i Uafgjorthedens Ligevægt, der bringer den paa Afstand. Og Hedenskabets høieste Mod var det, naar den Vise (hvis Alvor netop betegnedes ved ikke at forhaste sig med Forklaringen) formaaede at leve saaledes med Tanken om Døden: overvindende denne Tanke hvert Øieblik i sit Liv ved Uafgjortheden. Det jordiske Liv leves da ud, den Vise veed at Døden er til, han lever ikke tankeløst glemmende, at den er til, han mødes med den i Tanken, han gjør den afmægtig i Ubestemmeligheden, og denne er hans Seier over Døden; men Døden kommer ikke til at gjennemtrænge Livet omdannende.

Som det Uforklarlige kunde Døden synes at være den høieste Lykke. En saadan Forklaring forraader et Liv i Barnagtighed, Forklaringen er som dennes sidste Frugt: Overtro. Den Forklarende havde Barnets og Ynglingens Forestilling om det Behagelige og det Ubehagelige, og Livet gik hen, han saae sig bedragen, han blev ældre i Aar ikke i Sind, han greb intet Evigt: da samlede Barnagtigheden sig i ham til en overspændt Forestilling om at Døden skulde komme og lade Alt gaae i Opfyldelse; den blev nu den søgte Ven, den Elskede, den rige Velgjører, som havde alt Det at skjenke, hvad den Barnagtige forgjeves havde søgt opfyldt i Livet. Stundom tales der letsindigt og dumdristigt om denne Lykke, stundom veemodigt, stundom trænger den Forklarende sig endog høirøstet frem med sin Forklaring, og vil hjælpe Andre; men den forraader kun, hvorledes den Forklarende er i sit Indre, at han ikke fornam Alvorens Tilbagevirken, men er barnagtigt hastende fremad, barnagtigt haabende paa Døden, som han var det paa Livet.

Som det Uforklarlige kan Døden synes at være den største Ulykke. Men denne Forklaring angiver, at den Forklarende hænger feigt ved Livet, feigt maaskee ved dets Begunstigelse, feigt maaskee ved dets Lidelse, saa han frygter Livet, men frygter Døden endnu mere. Tilbagevirkende Kraft faaer Døden ikke, det vil sige ikke i Kraft af Opfattelsen, thi ellers virker den vel tilbage til at gjøre den Ene Lykkens Begunstigelse glædeløs, den Anden jordisk Lidelse haabløs.

Saa har Forklaringen ogsaa brugt andre betegnende Navne, den har kaldet Døden: en Overgang, en Forvandling, en Lidelse, en Strid, den sidste Strid, en Straf, Syndens Sold. Enhver af disse Forklaringer indeholder en heel Livs-Anskuelse. O, alvorlige Opfordring til den Forklarende! Let er det at sige dem alle udenad, let at forklare Døden, naar det ingen Overvindelse koster, ikke at ville forstaae, at Talen er om, at Forklaringen faaer tilbagevirkende Kraft i Livet. Hvorfor vil dog Nogen forvandle Døden til en Spot over sig; thi Døden trænger ikke til Forklaringen, den har vist aldrig anmodet nogen Tænker om at være sig behjælpelig. Men den Levende trænger til Forklaringen, og hvorfor? For at leve derefter.

Dersom Nogen saaledes mener, at Døden er en Forvandling, saa kan dette jo være ganske sandt, men sæt nu Dødens Uvished, der gaaer omkring ligesom Læreren og alle Øieblikke seer efter, om Discipelen er opmærksom, sæt nu den opdagede at den Forklarendes Mening var omtrent denne: jeg har et langt Liv for mig, tredive Aar ja maaskee fyrretyve, og saa kommer engang Døden som en Forvandling, hvad mon Læreren vilde tænke om denne Discipel, der end ikke havde fattet Uvishedens Bestemmelse ved Døden? Eller hvis Nogen mener, at det er en Forvandling som engang vil indtræde, og Dødens Uvished nu seer efter, og opdager, at han ikke ulig en Spiller forventer dette som en Begivenhed, der engang vil hænde, hvad mon Læreren vilde tænke om denne Discipel, der end ikke var opmærksom paa, at det i Dødens Afgjørelse er forbi, og at Forvandlingen ikke kan træde i Række med de øvrige Begivenheder som en ny Begivenhed, fordi der i Døden er sluttet af?

See, man kan have en Mening om fjerne Begivenheder, om en Naturgjenstand, om Naturen, om lærde Skrifter, om et andet Menneske, og saaledes om meget Andet, og naar man yttrer denne Mening, kan den Vise afgjøre, om den er rigtig eller urigtig. Derimod uleiliger Ingen den Menende med at betragte Sandhedens anden Side, om man nu virkelig har Meningen, om den ikke er Noget, man fremsiger. Og dog er denne anden Side ligesaa vigtig, thi Den er jo ikke blot afsindig, der siger det Meningsløse, men lige saa fuldt Den, der siger en rigtig Mening, naar denne slet og aldeles ingen Betydning har for ham. Det ene Menneske viser det andet den Tillid, den Anerkjendelse at antage, at det er hans Mening, naar han siger det. Ak, og dog er det saa let, saa saare let, at faae en sand Mening, ak, og dog er det saa vanskeligt, saa saare vanskeligt, at have en Mening, og at have den i Sandhed. Da nu Døden er Alvorens Gjenstand, saa er Alvoren atter her: at Døden betræffende skal man just ikke forhaste sig med at faae en Mening. Dødens Uvished tager sig jo bestandigt i al Alvor den Frihed at eftersee, om den Menende virkelig har denne Mening, det er, om hans Liv udtrykker den. I Forhold til Andet kan man yttre en Mening, og naar saa det fordres at man skal handle i Kraft af denne Mening, det er, vise at man har den, saa ere utallige Udflugter mulige. Men Dødens Uvished er Lærlingens strænge Hører; og naar han da fremsiger Forklaringen, saa siger Uvisheden til ham: nu, jeg skal undersøge om det er Din Mening, thi nu, nu i dette Øieblik er det forbi, for Dig, forbi, der er ikke Tanke om Udflugt, ikke et Bogstav at føie til, saa faaer jeg at see, om Du virkeligen meente hvad Du sagde om mig. Ak, al tom Forklaren, og al Ordgyden, og al Udpynten, og al Sammenkjeden af tidligere Forklaringer for at finde en endnu sindrigere, og al Beundring derfor og al Møie dermed: alt Dette er kun Adspredelse og Aandsfraværelse i Tankefjernhed – hvad mon Dødens Uvished tænker derom?

Derfor skal Talen afholde sig fra enhver Forklaring; som Døden er det Sidste af Alt, saa skal dette være det Sidste, der siges om den: den er uforklarlig. Uforklarligheden er Grændsen, og Udsagnets Betydning kun at give Dødens Tanke tilbagevirkende Kraft, gjøre den fremskyndende i Livet, fordi det med Dødens Afgjørelse er forbi, og fordi Dødens Uvished seer efter hvert Øieblik. Uforklarligheden er derfor ikke en Opfordring til at gjette Gaader, en Indbydelse til at være sindrig, men Dødens alvorlige Formaning til den Levende: jeg behøver ingen Forklaring, betænk Du, at det i denne Afgjørelse er forbi, og at den hvert Øieblik kan være forhaanden; see, dette er vel værd for Dig at betænke.

\bigskip
\begin{center}
$*$\quad$*$\\
$*$
\end{center}
\bigskip

M. T., maaskee synes det Dig, at Du af denne Tale kun faaer Lidet at vide; Du veed maaskee meget Mere selv, og dog skal den ikke have været forgjeves, hvis den betræffende Forestillingen om Dødens Afgjørelse har været Anledning til, at Du mindede Dig selv om, at det at vide Meget ikke er et ubetinget Gode. Maaskee synes det Dig, at Tanken om Døden kun er bleven forfærdende, og at den dog ogsaa har en mildere, en venligere Side for Betragtningen, at den trætte Arbeiders Længsel efter Hvile, den mødige Vandrers Hasten efter Udgangen, den Bekymredes Fortrøstning til Dødens smertestillende Søvn, den Misforstaaedes veemodige Trang til at slumre i Fred ogsaa er en skjøn og en berettiget Forklaring af Døden. Unægteligt! Men den læres ikke udenad, den læres ikke ved at læse om den, den erhverves langsomt, og først vel erhvervet af Den, som arbeidede sig træt i den gode Gjerning, som vandrede sig mødig paa den rette Vei, som bar Bekymringen for en retfærdig Sag, som misforstodes i en ædel Stræben, først saaledes vel erhvervet er den paa sit rette Sted, og en berettiget Tale i den Høiærværdiges Mund. Men den Unge tør ikke tale saaledes, at ikke den skjønne Forklaring, ligesom det vise Ord i en Daares Mund, i hans Mund skal blive en Usandhed. Og dette har jeg vel hørt, at Barnets og Ynglingens alvorlige Lærer i en sildigere Tid blev den Ældres og Modneres Ven, men jeg har aldrig hørt, idetmindste ikke af Nogen, af hvem jeg ønskede at lære, at det begyndte med, at Læreren strax blev Legebroder eller Barnet en Gammelmand, ei heller at da hiint Venskabs-Forhold i Sandhed indtraadte. Saaledes med Tanken om Døden. Har den ikke engang med Forfærdelse standset den Unges Liv, og kun brugt Alvoren til at holde Maade med Forfærdelsen, har Dødens Uvished ikke haft sin Underviisningens Tid, hvor den med Alvorens Strenghed oplærte ham: da har jeg aldrig hørt, idetmindste ikke af Nogen, i hvis Viden jeg kunde ønske Lod og Deel, jeg har aldrig hørt af nogen Saadan, at det da var Sandhed om Nogen kaldte Døden sin Ven, da han aldrig havde haft Andet i den end en Legebroder, naar han allerede i Ungdommen træt af Livet svigagtigt for at bedrage Livet talte om Dødens Venskab, naar han uden at have nyttet Livet som Olding svigagtigt for at bedrage sig selv talte om Dødens Venskab. – Den som her har talt, han er ung, endnu i den Lærendes Alder; han fatter kun Underviisningens Vanskelighed og Strenghed, o, at det maatte lykkes ham at gjøre det saaledes, at han netop derved blev værdig til engang at turde glæde sig ved Lærerens Venskab! Den som her har talet, han er jo ikke Din Lærer, m. T., han lader Dig jo blot, ligesom han selv er det, være Vidne til, hvorledes et Menneske søger at lære Noget af Tanken om Døden, hiin Alvorens Læremester, der ved Fødselen blev Enhver beskikket til Lærer for hele Livet, og som i Uvisheden altid er rede til at begynde Underviisningen, naar den forlanges. Thi Døden kommer ikke, fordi Nogen kalder paa den (det var kun Spøg, at saaledes den Svagere bød over den Stærkere), men saasnart Nogen aabner Uvisheden Indgang, saa er Læreren der. Læreren, som engang kommer for at holde Prøve og overhøre Discipelen, hvad enten han har villet nytte hans Underviisning eller ikke. Og denne Dødens Prøve, eller for med et mere brugt fremmed Ord at betegne det Samme, denne Livets sidste Examen er lige svær for Alle. Det er ikke som ellers, at den lykkeligt Begavede har let ved at bestaae, den ringe Begavede svært, nei, Døden afpasser Prøven i Forhold til Evnen, o, saa nøiagtigt, og Prøven bliver lige svær, fordi det er Alvorens Prøve.
\end{document}
